\documentclass[11pt]{article}
\usepackage[margin=1.05in]{geometry}
\usepackage{amsmath,amsthm,amssymb,mathtools,bm}
\usepackage{booktabs,array,enumitem}\usepackage[expansion=false]{microtype}
\usepackage{xcolor}
\usepackage[colorlinks=true,linkcolor=blue!60!black,citecolor=blue!60!black,urlcolor=blue!60!black]{hyperref}
\theoremstyle{plain}
\newtheorem{theorem}{Theorem}[section]\newtheorem{proposition}[theorem]{Proposition}\newtheorem{lemma}[theorem]{Lemma}
\newtheorem{corollary}[theorem]{Corollary}\newtheorem{conjecture}[theorem]{Conjecture}
\theoremstyle{definition}\newtheorem{definition}[theorem]{Definition}\newtheorem{remark}[theorem]{Remark}
\newtheorem{example}[theorem]{Example}
\newcommand{\R}{\mathbb R}\newcommand{\Z}{\mathbb Z}\newcommand{\E}{\mathbb E}\newcommand{\Var}{\operatorname{Var}}\newcommand{\Cov}{\operatorname{Cov}}
\newcommand{\relu}{\operatorname{relu}}\newcommand{\diag}{\operatorname{diag}}\newcommand{\tr}{\operatorname{tr}}
\newcommand{\Tr}{\operatorname{Tr}}\newcommand{\1}{\mathbf 1}\newcommand{\cA}{\mathcal A}\newcommand{\cB}{\mathcal B}
\newcommand{\cF}{\mathcal F}\newcommand{\cT}{\mathcal T}\newcommand{\cH}{\mathcal H}\newcommand{\aff}{\mathfrak{aff}}
\newcommand{\He}{\mathrm{He}}\newcommand{\Fa}{\Phi_{\mathrm F}}\newcommand{\cZ}{\mathcal Z}\newcommand{\Sym}{\operatorname{Sym}}
\newcommand{\Aff}{\mathrm{Aff}}\newcommand{\Ryshkov}{\mathcal R}
\title{\textsc{Localization in the commutant, stage 14:}\\ \large lattices, leaves and ellipsoids.\\
Klartag's process as a frozen-gap Dirac evolution, walls as sliding lattice points,\\ the resolution ellipsoid of a spectral truncation,
and the modular dilation of the non-commutative strata}
\author{}
\date{10 October 2026}
\begin{document}
\maketitle
\begin{abstract}
We take seriously the suggestion that the leaves of the network's wall foliation play the role of Klartag's lattice. We develop the
lattice side, the leaf side and the dictionary between them.
\begin{itemize}[leftmargin=*]
\item \emph{Klartag's process is spectral.} An ellipsoid free of the points of $L$ is the same as a flat torus $\R^n/L^*$ whose Laplacian
has spectral gap at least $4\pi^2$ in the corresponding metric. The contact points are the eigenfunctions at the gap: the first spectral
truncation of the Dirac operator. Klartag's constrained Dyson motion evolves the metric while freezing the gap, and it stops at a
Voronoi-perfect metric, one that its first spectral truncation determines. Local optima are eutactic, which is John's decomposition.
\item \emph{A leaf is a sliding lattice point (proved).} In the probe's precision coordinates, a wall's tangency constraint is exactly
Klartag's contact constraint $\langle y\otimes y,dA\rangle=0$, where $y$ is the wall's point of tangency.
\begin{itemize}
\item A wall is a continuum of lattice points, and only the tangency point binds.
\item Klartag's polyhedron becomes the intersection of infinitely many half-spaces, a convex matrix-fractional set.
\item The layer map $z=Wx$ is Klartag's normalizing transformation. It sends cells to orthants and the input ball to the state
ellipsoid, which is the frame the chain already works in.
\end{itemize}
\item \emph{The dilation can never be frozen.} For bias-free networks it is a free direction of every contact configuration (proved),
and biases remove it. This is the lattice-side reason for the critical gain mode of stage 13.
\item \emph{Spectral truncation is a resolution ellipsoid (proved, checked).} A kink at Mahalanobis distance $r$ has $k$-th chaos energy
$\He_{k-2}(r)^2\varphi(r)^2/k!$.
\begin{itemize}
\item Only $5$--$7\%$ of that energy lies below the Hermite turning point $k=r^2/4$, and essentially none below $r^2/8$.
\item Beyond the turning point the energy decays like $k^{-5/2}$; at $r=0$ the tail decays like $K^{-3/2}$, measured slope $-1.50$.
\item So truncation at chaos order $K$ is the ellipsoid of Mahalanobis radius $2\sqrt K$. Walls outside it are invisible to exponential
accuracy, and walls inside leave an algebraic tail.
\end{itemize}
\item \emph{The leaf space and its transverse measures.} The Gaussian restricts to transverse measures on every stratum, with the
mean as their codimension-one pairing (stage 11). The non-commutative strata are the bent crossings, whose isotropy is $\aff(1)$.
\begin{itemize}
\item $\aff(1)$ is isomorphic to the location--scale algebra of stage 12's hyperbolic upper tier, with the Euler (dilation) generator
matching the scale.
\item Its group $\Aff_+(1)$ is not unimodular, so the natural transverse weight is not a trace. Its modular flow is generated by the
logarithm of the dilation (Takesaki--Haagerup).
\item So the critical gain is the modular Hamiltonian of the network's non-commutative strata.
\item On the commutative part, the trace on $K_0$ of the tile groupoid takes values in the group generated by the cell masses: the
network's gap labels.
\end{itemize}
\item \emph{The code a local probe reveals.} Run Klartag's sticky evolution with walls as obstacles; its It\^o volume accounting checks
to $0.3\sigma$. At a maximal resolved probe, the whitened contact normals form a tight frame (John; checked to $10^{-14}$). The conic
Steiner formula holds in a network cell (checked); its face-dimension law is the global gate-count law. The rows' Gaussian
statistics play Siegel's role.
\end{itemize}
We end with what this changes for the chain: why the dilation must be carried, an adaptive resolution rule, positivity as a sticky
constraint, and the predictions. The checks are in \texttt{scripts/verify\_stage14.py}.
\end{abstract}
\tableofcontents
\section{The picture in one page}
\label{sec:summary}
\begin{center}\small
\begin{tabular}{>{\raggedright\arraybackslash}p{4.6cm}>{\raggedright\arraybackslash}p{9.6cm}}\toprule
Klartag (lattice) & the network (leaves)\\\midrule
lattice point $x\in L$ & a wall, acting through its sliding tangency point $x_*$ (Prop.~\ref{prop:sliding})\\
Ryshkov polyhedron $\{A:\ Ax\cdot x\ge1\}$ & curved Ryshkov body $\{A:\ w_a^\top A^{-1}w_a\le c_a\}$; a bundle of cones $\langle w\otimes w,mm^\top-\rho^2\Sigma\rangle\ge0$ over the cells\\
spectral gap of the flat torus & the chaos truncation of the probe's Ornstein--Uhlenbeck operator; resolution radius $2\sqrt K$ (Cor.~\ref{cor:resell})\\
contacts = gap eigenfunctions & walls at the resolution radius\\
frozen-gap Dyson motion of the metric & the sticky localization (It\^o accounting checked)\\
perfect form; Voronoi eutaxy & a perfect probe; John's decomposition, so contact normals form a tight frame (checked)\\
termination & never, for bias-free walls: the dilation is always free (Prop.~\ref{prop:dilfree})\\
normalizing map $T$ & the layer map $z=Wx$: cells to orthants, the ball to the state ellipsoid (the chain's frame)\\
Siegel's mean value & stage 12's field law; theta series $n(1+st)^{-1/2}$; the wall density at $s=1$\\
Rogers' second moment & correlations through the shared state; the leading one is the dilation (Mattis) mode\\
torus $\R^n/L$, Haar & leaf space of the wall foliation; Gaussian transverse measures on every stratum\\
$K_0$ and its trace & gap labels $\sum_C\Z\gamma(C)$ (cell masses)\\
space of lattices with scaling (Bost--Connes) & non-commutative strata $C^*(\Aff_+(1))$ with the dilation as modular flow (Thm~\ref{thm:modular})\\\bottomrule
\end{tabular}
\end{center}
\paragraph{The three new statements.}
\begin{itemize}[leftmargin=*]
\item \textbf{Leaves are sliding lattice points.} A wall imposes Klartag's constraint at every one of its points, and only the tangency
point binds. The resolved set is the intersection of the Ryshkov polyhedra of all leaf points.
\item \textbf{Spectral truncation is a resolution ellipsoid.} The kink at Mahalanobis distance $r$ puts only $5$--$7\%$ of its chaos energy below the
turning point $r^2/4$ (essentially none below $r^2/8$), then decays as $k^{-5/2}$. Order $K$ resolves exactly the walls inside radius $2\sqrt K$, with tail $\varphi(r)K^{-3/2}/3\pi$.
\item \textbf{The dilation is modular.} At bent crossings the isotropy $\aff(1)$ is the upper tier's location--scale algebra. Its group is not
unimodular, and the modular flow of its Plancherel weight is the dilation. The critical gain is the modular Hamiltonian of the
non-commutative strata, the free direction of every contact configuration, and the conserved charge of the layer map.
\end{itemize}
\paragraph{For the chain.}
\begin{enumerate}[leftmargin=*]
\item Carry the dilation (it cannot be pinned).
\item Resolution follows $r^2/4$, so effort follows the wall density.
\item Positivity on the truncated operator system is a truth-free sticky constraint.
\item The cell-adapted probe and its frame code are a local model, not yet an estimator.
\end{enumerate}

\section{Klartag's process is a spectral process}
\label{sec:klartag}
\subsection{Lattice-free ellipsoids are flat tori with a spectral gap}
Let $L\subset\R^n$ be a lattice, $A$ positive definite, and $E_A=\{x:\ Ax\cdot x<1\}$. Consider the flat torus $T_A=\R^n/L^*$ with the
metric whose inverse is $A$ ($\langle u,v\rangle=u^\top A^{-1}v$). Its Laplace eigenfunctions are the characters
$e_\xi(x)=e^{2\pi i\langle\xi,x\rangle}$, $\xi\in L$, with eigenvalues $4\pi^2A\xi\cdot\xi$, and its volume is $\mathrm{covol}(L)^{-1}\det A^{-1/2}$.
\begin{proposition}[known; restated]\label{prop:torus}
The following hold.
\begin{itemize}[leftmargin=*]
\item $E_A\cap L=\{0\}$ if and only if $\lambda_1(T_A)\ge4\pi^2$, and $\mathrm{vol}(E_A)=\kappa_n\,\mathrm{covol}(L)\,\mathrm{vol}(T_A)$. So Klartag's theorem
says that some flat torus has spectral gap $4\pi^2$ and volume $\gtrsim n^2/(\kappa_n\,\mathrm{covol}L)$: the packing problem is the
spectral-gap problem for flat tori.
\item The \emph{contacts} $\partial E_A\cap L$ are exactly the frequencies of the eigenfunctions at the gap. The first spectral projection
$P=\1_{[0,4\pi^2]}(\Delta_A)$ spans the constants and those characters. Connes--van Suijlekom's truncated operator system
$PC(T_A)P$ is spanned by the characters $e_{\xi-\eta}$ with $\xi,\eta$ contacts: the \emph{difference set} of the contact code.
\end{itemize}
\end{proposition}

\subsection{The process as a frozen-gap metric evolution}
Klartag's $dA_t=\pi_t(dW_t)$ is a Dyson Brownian motion of the metric, projected so that $d\langle A\xi,\xi\rangle=\langle\xi\otimes\xi,dA\rangle=0$
for every contact. In spectral words, the metric diffuses while every eigenvalue that reaches the gap stays frozen there.
\begin{itemize}[leftmargin=*]
\item \textbf{Volume.} By It\^o, $d\log\det A_t=\langle A_t^{-1},\pi_tdW_t\rangle-\frac12\sum_{i,j}|\pi_t(u_i\otimes_su_j)|^2/(\lambda_i\lambda_j)\,dt$. The torus volume
therefore grows in expectation at a rate set by the number of \emph{unfrozen} metric directions,
$\dim F_t=\frac{n(n+1)}2-\mathrm{rank}\{\xi\otimes\xi\}$.
\item \textbf{Termination.} The process stops when $F_t=\{0\}$, i.e.\ when the contact outer products span $\mathrm{Sym}_n$. This is
Voronoi's \emph{perfect} form. In spectral terms, the first spectral truncation determines the metric: \emph{rigidity at finite
resolution}.
\item \textbf{Optimality.} At a local maximum of volume, the KKT conditions give $A^{-1}=\sum\lambda_\xi\,\xi\otimes\xi$ with $\lambda_\xi>0$. This is
Voronoi's \emph{eutaxy}, which in John position reads $\sum\lambda_\xi u_\xi\otimes u_\xi=\mathrm{Id}$: a tight frame. Voronoi's theorem
says the locally densest lattices are exactly the perfect and eutactic ones.
\item \textbf{Compactness.} By Mahler's criterion, tori with gap at least $4\pi^2$ and bounded volume form a compact set. The process
lives in the compact thick part of the space of lattices, and the contacts are its boundary.
\end{itemize}
\emph{Check.} $A_2$ has $6$ minimal vectors whose outer products have rank $3=\dim\mathrm{Sym}_2$ (perfect), with $\sum x\otimes x=3\,\mathrm{Id}$
(eutactic). $\Z^2$ has $4$ minimal vectors of rank $2$, so it is not perfect.

\paragraph{Why this matters for us.} Stated spectrally, Klartag's mechanism has four parts: a metric (Dirac) evolution; a resolution
threshold (the gap); objects at the threshold that freeze (contacts); and a volume accounting by unfrozen directions. Our setting has
all four. The metric is the localization probe's covariance, the threshold is the chaos truncation (Section~\ref{sec:resolution}), and
the objects at the threshold are walls (Section~\ref{sec:leaves}).

\section{Leaves as sliding lattice points}
\label{sec:leaves}
\subsection{A wall is a continuum of lattice points}
A \emph{probe} is a Gaussian localization state, written as an ellipsoid $E(m,A)=\{x:\ A(x-m)\cdot(x-m)<1\}$. Here $A$ is the precision at the
chosen Mahalanobis radius. The probe is \emph{resolved} for a wall $H=\{w\cdot x=b\}$ if $E\cap H=\emptyset$, and in \emph{contact} if $H$ is
tangent to $\partial E$.
\begin{proposition}[the leaf as a sliding lattice point; proved]\label{prop:sliding}
\begin{enumerate}[leftmargin=*]
\item $E\cap H=\emptyset$ if and only if $\langle(x-m)\otimes(x-m),A\rangle\ge1$ for every $x\in H$. The wall therefore imposes Klartag's
lattice-point constraint once for each of its points. The binding point is the tangency point
$x_*=m+A^{-1}w\,\frac{b-w\cdot m}{w^\top A^{-1}w}$, and the inequality collapses to $w^\top A^{-1}w\le(w\cdot m-b)^2=:c$.
\item At a contact, with fixed centre, the infinitesimal constraint is $\langle(x_*-m)\otimes(x_*-m),\,dA\rangle=0$. This is Klartag's contact
constraint, with the lattice point replaced by the wall's tangency point, which slides as $A$ moves.
\item The resolved set $\Ryshkov(m)=\{A:\ w_a^\top A^{-1}w_a\le c_a\ \forall a\}$ is convex (matrix-fractional). It is the intersection of the
Ryshkov polyhedra of all points of all leaves: a \emph{curved Ryshkov body}.
\end{enumerate}
\end{proposition}
\begin{proof}
The $A$-distance from $m$ to $H$ is $|w\cdot m-b|/\sqrt{w^\top A^{-1}w}$, attained at $x_*$. Differentiating, $d(w^\top A^{-1}w)=-\langle A^{-1}ww^\top A^{-1},dA\rangle$, and
$A^{-1}w\propto x_*-m$. Finally, $w^\top A^{-1}w$ is convex in $A\succ0$.
\end{proof}
This is the precise sense in which leaves play the lattice. A lattice point is a leaf of dimension $0$. A wall is a leaf of dimension
$n-1$ that acts through one moving point, its point of tangency. Klartag's polyhedron, cut by finitely many fixed points near the
origin, becomes a convex body cut by a family of sliding points.

\subsection{With the centre free: a cone in the space of states}
In terms of the covariance $\Sigma$ at Mahalanobis radius $\rho$, a probe $(m,\Sigma)$ is resolved for $H$ iff
$\langle w\otimes w,\ mm^\top-\rho^2\Sigma\rangle-2b\,w\cdot m+b^2\ge0$. For the bias-free first layer this is
\[
\big\langle w_a\otimes w_a,\ M\big\rangle\ \ge0\quad\forall a,\qquad M=mm^\top-\rho^2\Sigma .
\]
This is a polyhedral cone in the variable $M$, and the rank-one matrices $w_a\otimes w_a$, the rows of $W_1$, play the lattice vectors. At
depth $l$ the walls are piecewise linear. Inside a cell, $w_a$ is replaced by the gated Jacobian row $\nabla z_{l,a}$ (and $b$ by $0$, by
homogeneity). The resolved set is therefore a \emph{bundle} of such cones over the cells: the lattice varies along the leaf space of the
tile groupoid.

\subsection{The layer map is Klartag's normalizing transformation}
Klartag ends by applying $T$ with $T(E)=$ ball, which turns the lattice into a packing. Our network applies the dual move at every
layer:
\begin{itemize}[leftmargin=*]
\item $z=W_1x$ sends the walls to the coordinate hyperplanes, the cells to the $2^n$ orthants, and the input ball to the ellipsoid
$N(0,W_1W_1^\top)$;
\item in $z_l$-coordinates the layer-$l$ walls are again coordinate hyperplanes, and the state is the closure's $N(\mu_l,C_l)$.
\end{itemize}
So the chain lives in the frame where the ``lattice'' is fixed and standard (the coordinate arrangement, whose face lattice is the
Boolean lattice of gate patterns) and the ellipsoid moves. The fields $r_a=\mu_a/\sqrt{C_{aa}}$ are the Mahalanobis distances to the walls,
and contacts are the neurons with $|r_a|=\rho$.

\subsection{The dilation can never be frozen}
\begin{proposition}[proved]\label{prop:dilfree}
For bias-free walls, the dilation $(\delta m,\delta\Sigma)=(m,2\Sigma)$ preserves every resolution inequality and every tangency, because each
constraint is homogeneous of degree $2$ in $(m,\Sigma^{1/2})$. It therefore lies in the free subspace $F$ of every contact configuration.
More generally, every affine vector field that preserves the contact walls acts inside $F$. With biases ($b\ne0$) the dilation is not
free.
\end{proposition}
\emph{Check:} along the dilation, $dg_a-2g_a=0$ exactly for homogeneous walls, and is $O(1)$ for affine walls.
\begin{corollary}[the lattice-side reason for the critical gain]
With homogeneous walls, Klartag's process never terminates: the probe keeps diffusing along the log-scale, whatever contacts it has
collected. No local (contact) information pins the scale. The chain must therefore carry it as a conserved charge, which is stage 13's
dilation package.
\end{corollary}
\begin{conjecture}[falsifiable]
In networks with biases, the gain mode's error-transport factor is strictly below $1$, with a gap that grows with the bias scale.
Biases are lattice points that break the dilation.
\end{conjecture}

\subsection{Depth sheds contacts}
Klartag's ellipsoid accumulates contacts. The network's state sheds them. As depth grows, the fields spread as $N(0,t_l)$ (stage 12),
so the walls leave the resolution ellipsoid. The number still inside radius $\rho$ is $n\operatorname{erf}(\rho/\sqrt{2t_l})\approx n\rho\sqrt{2/\pi t_l}$, the
wall-density law. Cooling is Klartag's process run backwards.

\section{Spectral truncation is a resolution ellipsoid}
\label{sec:resolution}
\subsection{The kink chaos law}
For a neuron whose pre-activation is $z=\mu+s\xi$ (field $r=\mu/s$), expand $\relu(z)/s$ in the Wiener chaos of $\xi$. The first chaos
carries $\Phi(r)^2$. The rest is the kink.
\begin{theorem}[the kink chaos law; proved, and checked]\label{thm:kink}
The $k$-th chaos energy of $\relu(r+\xi)$ is, for $k\ge2$,
\[
E_k(r)=\frac{\He_{k-2}(r)^2\,\varphi(r)^2}{k!},\qquad \sum_{k\ge2}E_k(r)=\Var\relu(r+\xi)-\Phi(r)^2 .
\]
By the Plancherel--Rotach asymptotics:
\begin{itemize}[leftmargin=*]
\item $E_k(r)$ is exponentially small for $k<r^2/4$, where $r$ lies outside the oscillatory region $|r|<2\sqrt{k}$ of $\He_{k-2}$;
\item beyond the turning point, averaging over the oscillation, $E_k(r)\approx\dfrac{\varphi(r)}{\pi k^2\sqrt{4k-r^2}}$;
\item hence the tail beyond order $K$ is $\sum_{k>K}E_k(r)\approx\dfrac{\varphi(r)}{3\pi}K^{-3/2}$ for $K\gg r^2/4$.
\end{itemize}
\end{theorem}
\begin{proof}
Stein: $\E[\relu(r+\xi)\He_k(\xi)]=\E[\relu^{(k)}(r+\xi)]=\E[\delta^{(k-2)}(r+\xi)]=\pm\He_{k-2}(r)\varphi(r)$. The energy is the square over $k!$. The asymptotics are
those of the normalized Hermite functions $\psi_j(x)=\He_j(x)e^{-x^2/4}/((2\pi)^{1/4}\sqrt{j!})$. Since $E_k=\psi_{k-2}(r)^2\varphi(r)/(k(k-1))$, and
$\psi_j^2$ is exponentially small outside $|x|<2\sqrt j$ with oscillation-averaged value $1/(\pi\sqrt{4j-x^2})$ inside (WKB; the oscillator
$-\partial_x^2+x^2/4$ at energy $j+\frac12$), the asymptotics follow.
\end{proof}
Checked:
\begin{itemize}[leftmargin=*]
\item The sum rule holds to $10^{-7}$ relative.
\item The share of the kink energy below the turning point $k=r^2/4$ is $0.073$, $0.056$, $0.049$ at $r=4,8,12$; below $r^2/8$ it is
$0.003$, $0$, $0$.
\item The cumulative tail exponent beyond $4r^2$ is $-1.52$.
\item At $r=0$ the tail is $4.214\times10^{-5}$, $1.337\times10^{-6}$, $4.185\times10^{-8}$ at $K=10^2,10^3,10^4$ (slopes $-1.50$, $-1.51$). The predicted
$\varphi(0)K^{-3/2}/3\pi$ gives $1.339\times10^{-6}$ at $K=10^3$.
\end{itemize}

\subsection{Truncation at order $K$ is the ellipsoid of radius $2\sqrt K$}
The Ornstein--Uhlenbeck operator of the probe is the number operator: the $k$-th chaos is its eigenspace with eigenvalue $k$. It is
the Dirac-type operator whose spectral truncation at level $K$ is what a chain truncated at chaos order $K$ keeps.
Theorem~\ref{thm:kink} is its Weyl correspondence for kinks.
\begin{corollary}[resolution ellipsoid; proved]\label{cor:resell}
Truncating at chaos order $K$ is, up to exponentially small errors, the probe ellipsoid of Mahalanobis radius $2\sqrt K$.
\begin{itemize}[leftmargin=*]
\item A wall outside it ($r^2>4K$) leaves essentially nothing in the retained chaos. Its whole nonlinear content is at most $\varphi(r)/r^3$.
\item A wall inside it is captured up to an algebraic tail $\propto\varphi(r)K^{-3/2}$.
\end{itemize}
The chaos truncation is Klartag's resolution threshold. The walls inside the $2\sqrt K$-ellipsoid are the ``lattice points in the
ellipsoid'', and those near its boundary are the contacts.
\end{corollary}
\begin{corollary}[a truncation-error law; derived, heuristic in its use]
With stage 12's field law ($r\sim N(0,t_l)$ across neurons), the chaos energy neglected at order $K$ in layer $l$ is
\[
\sum_a\sum_{k>K}E_k(r_a)\ \approx\ \frac{n\,K^{-3/2}}{3\pi}\,\E\varphi(r)=\frac{n\,K^{-3/2}}{3\pi}\cdot\frac{\sin(\theta_l/2)}{\sqrt\pi}.
\]
Raising the truncation from $K$ to $K+1$ removes the fraction $1-(K/(K+1))^{3/2}$ of it: $46\%$ ($2\to3$), $35\%$ ($3\to4$),
$28\%$ ($4\to5$). For comparison, v56's chaos grading (roughly a step from second to third chaos in the $\kappa_4$ diagonal) took $21\%$ off the raw MSE. These are the same
order, but the MSE has other sources.
\end{corollary}
At $K=3$ the resolution radius is $2\sqrt3=3.46$. The fraction of neurons inside it is $\operatorname{erf}(\sqrt{2K/t_l})$: $0.92$ at $l=8$ ($t\approx4$) and
$0.68$ at $l=16$ ($t\approx12$). This matches stage 12's measurement that the marginal memory vanishes for $|r|\ge3$.

\subsection{The chain's state lives on a truncated operator system}
Truncating at order $K$ replaces the commutative algebra of the probe by the operator system $\mathcal E_K=P_K\,L^\infty(\gamma_{m,\Sigma})\,P_K$
(Connes--van Suijlekom). The chain's carried statistics are a functional on $\mathcal E_K$. For that functional to come from a probability
law it must be positive, a truncated moment condition. Positivity defines a convex body whose boundary points are functionals with a
kernel. These are the analogue of Klartag's contacts: directions in which the chain's pseudo-law has no room left. Section~\ref{sec:chain}
uses this as a truth-free sticky constraint.

\section{The leaf space, its transverse measures, and the modular dilation}
\label{sec:groupoid}
\subsection{Strata, leaves and transverse measures}
Stage 11 set up the wall foliation of input space at depth $l$. Its leaves are the open faces of the activation tiling: cells
(codimension $0$), walls ($1$), and crossings ($\ge2$). Its groupoid algebra has a finite composition series over the strata, with
subquotients $C^*(\text{isotropy})\otimes\mathcal K$:
\begin{itemize}[leftmargin=*]
\item $C_0(\R^j)$ at normal crossings, where the isotropy is abelian;
\item $C^*(\Aff_+(1))$ at bent crossings, where a deep wall meets a wall that feeds it.
\end{itemize}
The Gaussian $\gamma$ induces a transverse measure on every stratum by disintegration. On a stratum $S$ of codimension $j$ it is the
Gaussian surface density on $S$, with the normal space $\R^j$ as local transversal.
\begin{proposition}[the mean is leaf-space integration; stage 11, restated]\label{prop:leafint}
\[
\E f=(2\pi)^{-1/2}\sum_{\mathrm{codim}\,F=1}\kappa_F\,\alpha_F ,
\]
the codimension-one transverse measure paired with the kink cocycle $\kappa_F$.
The cumulant corrections of order $k$ pair the codimension-$j$ transverse measures, $j\le k-1$, with higher kink cocycles. Only bent
chains carry mixed terms (stage 12: Kikuchi level $=$ codimension). The mean, and every correction to the closure, is an integral over
the leaf space.
\end{proposition}

\subsection{The tile groupoid and the network's gap labels}
Let $R_l$ be the relation ``same cell at depth $l$''. It is proper, and $C^*(R_l)\cong\bigoplus_C\mathcal K(L^2(C))$, so $K_0(C^*(R_l))\cong\Z^{(\text{cells})}$.
\begin{proposition}[proved, with the trace normalized by $\gamma$]\label{prop:gap}
The trace induced by the Gaussian transverse measure gives each cell $C$ the weight $\gamma(C)$. Its range on $K_0$ is the
\emph{gap-labelling group}
\[
\mathrm{Gap}_l=\tau_*K_0(C^*(R_l))=\sum_C\Z\,\gamma(C)\ \subset\R .
\]
For every self-adjoint element $H$ of $C^*(R_l)$ (a cell-covariant observable, for instance any function of the gate code), the
integrated density of states $\tau(\1_{(-\infty,E]}(H))$ at a spectral gap lies in $\mathrm{Gap}_l$. Refinement with depth gives
$\mathrm{Gap}_l\subset\mathrm{Gap}_{l+1}$.
\end{proposition}
This is the commutative half of the leaf space. Its K-theoretic ``lattice'' is the module of cell masses, the global code of gate
patterns weighted by the Gaussian. Bellissard's gap labelling for quasicrystals has the same shape, with patch frequencies in place
of cell masses.

\subsection{The non-commutative strata and their modular dilation}
\begin{theorem}[structural; proved at the Lie-algebra level]\label{thm:affiso}
Stage 11 found the isotropy algebra at a bent crossing, $\aff(1)=\langle E,\vartheta\rangle$ with $[E,\vartheta]=\vartheta$, where $E$ is the Euler (dilation) field
and $\vartheta$ the bending shear. It is isomorphic to the Lie algebra of the location--scale group $\Aff_+(1)=\{x\mapsto\sigma x+\mu\}$, which acts
simply transitively on univariate Gaussians. That group is stage 12's hyperbolic upper tier. The isomorphism sends $E$ to the
scale generator $D$ and $\vartheta$ to the translation $P$ ($[D,P]=P$).
\end{theorem}
\begin{theorem}[the modular flow is the dilation; known ingredients, applied]\label{thm:modular}
$\Aff_+(1)$ is not unimodular. Its left Haar measure is $d\sigma\,d\mu/\sigma^2$, its right Haar measure $d\sigma\,d\mu/\sigma$, and its modular
function $\Delta(\sigma,\mu)=\sigma^{-1}$. The Plancherel weight $\varphi$ on the group von Neumann algebra $L(\Aff_+(1))$ is faithful, normal and
semifinite, but not a trace. Its modular group is (Takesaki, Haagerup)
\[
\sigma^\varphi_s\big(\lambda_{(\sigma,\mu)}\big)=\Delta(\sigma,\mu)^{is}\,\lambda_{(\sigma,\mu)}=\sigma^{-is}\,\lambda_{(\sigma,\mu)} .
\]
The modular Hamiltonian $\log\Delta=-\log\sigma$ is the logarithm of the dilation.
\end{theorem}
\begin{corollary}[the critical gain is modular]
On the non-commutative strata, the natural transverse weight is not a trace, and its modular flow is the dilation. The mode that
stage 13 found critical (transport factor $1$, a conserved charge) is the modular Hamiltonian of the network's non-commutative
strata. Three statements are one: the scale is conserved (homogeneity), never frozen by contacts (Proposition~\ref{prop:dilfree}), and
generates the modular flow where the foliation is non-commutative.
\end{corollary}
\begin{remark}[what follows, and what does not]
\begin{itemize}[leftmargin=*]
\item By Connes' Thom isomorphism, the group $C^*$-algebra of $\Aff_+(1)$ has $K_0=\Z$ and $K_1=0$. Gap labels on the bent strata cannot
come from a trace. They need the dilation-twisted (KMS) weight, that is, the twisted cyclic cohomology of Connes and Moscovici.
\item Tomita--Takesaki gives $JL(G)J=R(G)$, so the commutant of the left group algebra is the right one. Stage 11 had forward activations
and backward sensitivities as left- and right-harmonic states. Reading them as the left and right sides of the bent-stratum algebra is
an analogy we have not made precise.
\end{itemize}
\end{remark}

\subsection{An analogy with the space of lattices}
Klartag's process lives, in its dual form, on the space of unimodular lattices $X_n=SL_n(\Z)\backslash SL_n(\R)$, in its Mahler-compact thick
part. The Bost--Connes and Connes--Marcolli systems make a non-commutative space of lattices with the \emph{scaling} as time evolution.
Its KMS states encode the arithmetic, with partition function $\zeta(\beta)$. For the network, the same dilation is the modular flow of the
non-commutative strata. This suggests one well-defined object, the \emph{radial partition function}
\[
Z_l(\beta)=\frac{\E|h_l(x)|^\beta}{\E|x|^\beta}=\int_{S^{n-1}}|h_l(\hat x)|^\beta\,d\sigma(\hat x),
\]
the Mellin transform of the gain. Its logarithm generates the gain's cumulants: $\frac{d^2}{d\beta^2}\log Z|_{0}=\Var\log|h_l(\hat x)|$, the dilation
charge, about $V_l/4$ minus the input radius's share. It is the exact object a dilation-consistent closure must transport (stage 13).
The arithmetic content of Bost--Connes has no counterpart here, and we do not claim one.

\section{The ellipsoid process on the leaves, and the code it reveals}
\label{sec:process}
\subsection{The sticky localization}
\begin{definition}
Fix a centre $m$ and walls $\{w_a\cdot x=b_a\}$. Start from a resolved probe precision $A_0=a_0\mathrm{Id}$. Run Klartag's evolution on the precision,
\[
dA_t=\sigma\,\pi_t(dW_t),\qquad \pi_t=\text{projection onto }\{B:\ \langle y_a\otimes y_a,B\rangle=0\ \text{for every contact }a\},
\]
where $y_a=A^{-1}w_a$ is the direction to wall $a$'s tangency point (Proposition~\ref{prop:sliding}). A wall becomes a contact when it first
touches the probe, and stays one. The process freezes when the $y_a\otimes y_a$ span $\mathrm{Sym}_n$.
\end{definition}
By It\^o, $\E\log\det A_T=\log\det A_0-\frac{\sigma^2}2\E\int_0^T\delta_tdt$, with $\delta_t=\sum_{i,j}|\pi_t(u_i\otimes_su_j)|^2/(\lambda_i\lambda_j)$. Since $\log\det$ is
concave, the expected probe volume grows at a rate set by the free directions.
\emph{Check} ($n=3$, $14$ affine walls, $200$ paths, $T=2$):
\begin{itemize}[leftmargin=*]
\item $\E\log\det A_T=18.700\pm0.050$, against the It\^o prediction $18.683\pm0.010$, from $\log\det A_0=19.114$;
\item contacts accumulate to $1.9$ on average and up to $6=\dim\mathrm{Sym}_3$, i.e.\ perfect.
\end{itemize}
With homogeneous walls and the centre free, Proposition~\ref{prop:dilfree} keeps the dilation free forever. The probe then never
freezes, and its log-scale performs the free walk that is the critical mode.

\subsection{The maximal resolved probe and its code}
\begin{proposition}[Klartag's terminal state meets John; proved, and checked]\label{prop:john}
The locally maximal resolved probe about $m$ is the polar of the L\"owner ellipsoid of the dual wall points $\pm w_a/|w_a\cdot m-b_a|$. It touches
the contact walls at Mahalanobis distance $1$, and its precision satisfies the eutaxy (John) condition
\[
A=\sum_a\lambda_a\,x_a\otimes x_a,\qquad \lambda_a\ge0,\qquad \sum\lambda_a=n,\qquad x_a=\frac{w_a}{|w_a\cdot m-b_a|}.
\]
This is the polar form of Voronoi's eutaxy: for lattice points the inverse of the ellipsoid matrix is the frame sum, while for walls the
precision is. In the whitened frame $u_a=A^{-1/2}x_a$, $\sum_a\lambda_au_a\otimes u_a=\mathrm{Id}$: the contact normals form a weighted tight frame, a
spherical $2$-design.
\end{proposition}
\emph{Check} ($n=4$, $9$ walls): Khachiyan's algorithm gives $4$ active walls, all at Mahalanobis distance $1.0000$, with
$\|\sum\lambda_au_a\otimes u_a-\mathrm{Id}\|=1.3\times10^{-14}$ and $\sum\lambda_a=4.0000$.

This is the transformation that turns a local process into a code. Klartag's $T$ maps his terminal ellipsoid to the ball and his
lattice to a packing. Our whitening $A^{1/2}$ (the map $x\mapsto A^{1/2}(x-m)$) maps the maximal resolved probe at an input to the unit ball, and the walls around that
input to a tight frame of contact normals. The local geometry of linearity around an input, expressed in its own natural metric, is a
spherical design: the neighbourhood of every input carries a small frame code.

\subsection{Cells as cones: the conic Steiner formula and the global gate-count law}
\textsc{ReLU} is the metric projection onto the orthant, $\relu(z)=\Pi_{\R^n_+}(z)$. Moreau's decomposition $z=\Pi_K z+\Pi_{K^\circ}z$, i.e.\
$\relu(z)-\relu(-z)=z$, is the cone--polar duality, our analogue of lattice--dual lattice. For any closed convex cone $K$ and
$g\sim N(0,\mathrm{Id})$, the conic Steiner formula (McCoy--Tropp) says
\[
\E\,f\big(|\Pi_Kg|^2,|\Pi_{K^\circ}g|^2\big)=\sum_kv_k(K)\,\E f(\chi^2_k,\chi^2_{n-k}),
\]
where $v_k$ is the law of the dimension of the face that $\Pi_Kg$ lands in (the conic intrinsic volumes).
\begin{itemize}[leftmargin=*]
\item For the orthant, $v_k$ is binomial: the face is the gate pattern. For orthogonal weights the first layer is exactly this, so
$v_k$ is the global gate-count law.
\item \emph{Check}, in a layer-1 cell $C=W^{-1}\R^n_+$ with Gaussian $W$ ($n=10$, $2\times10^4$ samples, exact projections by NNLS):
\begin{itemize}
\item $\E|\Pi g|^2=3.179\pm0.020$ against $\sum kv_k=3.188$;
\item $\E|\Pi g|^4=17.89\pm0.24$ against $17.93$;
\item $\E|\Pi g|^2|\Pi^\circ g|^2=20.12\pm0.16$ against $20.33$.
\end{itemize}
\item The cell's statistical dimension is $3.19$, against $n/2=5$ for orthogonal weights. Gaussian weights make the cell thinner than an
orthant.
\end{itemize}

\subsection{Siegel's principle for the network's rows}
Klartag's analysis rests on Siegel's mean value theorem: a random lattice is, on average, Lebesgue measure. Our analogue is stage
12's self-localization. Through Gaussian rows independent of the state, the fields $r_a$ of a probe are, across rows, $N(0,t)$-distributed.
The ``theta series'' of the wall code relative to the probe therefore has the Siegel mean
\[
\Theta(s)=\sum_a e^{-sr_a^2/2},\qquad \E\,\Theta(s)=n\,(1+st)^{-1/2}.
\]
Its value at $s=1$ is $\sqrt{2\pi}$ times the wall density, $n/\sqrt{1+t}$, which is stage 12's wall-density law. Rogers' second-moment
formula has an analogue as well. Given the state, the rows are independent (Poisson-like), and the correlations come only through the
shared state. The leading one is the collective (dilation) mode, the Mattis component of stage 12, which is once more the charge.

\section{What it changes for the chain}
\label{sec:chain}
\subsection{What the dictionary says about the chain}
\begin{enumerate}[leftmargin=*]
\item \textbf{The dilation must be carried, because it cannot be pinned.}
\begin{itemize}
\item In a bias-free network the dilation is a free direction of every contact configuration (Proposition~\ref{prop:dilfree}) and the
modular flow of the non-commutative strata (Theorem~\ref{thm:modular}).
\item So no closure built from local (contact) data fixes it, and errors along it are not damped. That is note XLIV's critical mode,
now seen from three sides.
\item The remedy is stage 13's: carry the charge (the radial partition function's second cumulant) and write the variance, $\kappa_3$ and
$\kappa_4$ along the dilation as one package.
\item \emph{Prediction:} networks with biases show a damped gain mode.
\end{itemize}
\item \textbf{Resolution is set by $r^2/4$, so effort should follow the wall density.} By Corollary~\ref{cor:resell}, a neuron at field $r$ needs chaos
order about $r^2/4$ before its kink is seen at all, and contributes at most $\varphi(r)/r^3$.
\begin{itemize}
\item The chain's truncation error at order $K$ concentrates on neurons with $r^2\lesssim4K$, weighted by $\varphi(r)K^{-3/2}$.
\item \emph{Prediction (testable on note XLI's dumps):} the per-neuron benefit of v56's chaos correction is supported on $|r|\lesssim2\sqrt3$,
with a $\varphi(r)$-weighted profile.
\item \emph{Prediction:} the next chaos order removes about $35\%$ of the remaining chaos-truncation energy.
\end{itemize}
\item \textbf{Positivity is a sticky constraint.} The chain's state is a functional on the truncated operator system $\mathcal E_K$. If it
leaves the positive cone, the pseudo-law has negative mass somewhere near a wall.
\begin{itemize}
\item Klartag's rule is to project each layer's update onto the face it touches and continue there.
\item This is truth-free and costs $O(n^2)$ per layer for the marginal (one-neuron) moment matrices.
\item A first diagnostic: count how often v56's per-neuron moment matrices (orders $\le4$) fail to be positive semidefinite.
\end{itemize}
\item \textbf{The cell-adapted Gaussian.} The maximal resolved probe at an input is the natural local Gaussian of its cell, and its frame
code records the cell's local wall geometry. It is a candidate local model for future hybrid estimators. Sampling is roughly
$1000\times$ too expensive at this budget (stage 13), so this is no algorithm yet.
\end{enumerate}

\subsection{Open problems}
\begin{enumerate}[leftmargin=*]
\item \textbf{Klartag's estimate for walls.} Prove that, for Gaussian rows, the sticky localization reaches a resolved probe of
log-volume $\ge c\,\cdot$(free dimensions) before contacts saturate. Klartag's $n^2$-versus-$n$ gain would be the ratio between the
best Gaussian probe and the best ball inside a typical cell.
\item \textbf{Twisted gap labels.} Compute the dilation-twisted (KMS) pairing of the bent strata, the non-commutative part of the
gap-labelling group, and decide whether it equals a correction term of the mean.
\item \textbf{Left and right.} Make precise whether the Tomita conjugation $JL(G)J=R(G)$ at bent strata exchanges stage 11's forward and
backward harmonic states.
\item \textbf{The radial partition function.} Find a recursion for $Z_l(\beta)$ through a layer. The dilation package of stage 13 uses
only its second cumulant.
\item \textbf{Codes across inputs.} The tight frames of Proposition~\ref{prop:john} vary with the input. Is their distribution over the
sphere a function of the internal clock alone, as stage 12's field law is?
\end{enumerate}

\section*{Verification}
The checks are run by \texttt{scripts/verify\_stage14.py}; its output is saved next to this note.
\begin{center}\small
\begin{tabular}{>{\raggedright\arraybackslash}p{5.0cm}>{\raggedright\arraybackslash}p{9.8cm}}\toprule
identity & result\\\midrule
1. Perfect and eutactic (Section~\ref{sec:klartag}) & $A_2$: $6$ minimal vectors, rank $3$, $\sum x\otimes x=3\,\mathrm{Id}$; $\Z^2$: rank $2$ (not perfect)\\
2. Kink chaos law (Thm~\ref{thm:kink}) & sum rule to $10^{-7}$; share below $r^2/4$: $0.073$, $0.056$, $0.049$ ($r=4,8,12$); below $r^2/8$: $\le0.003$; tail exponent $-1.52$; at $r=0$ tail
$1.337\times10^{-6}$ at $K=10^3$ against $\varphi(0)K^{-3/2}/3\pi=1.339\times10^{-6}$, slope $-1.50$\\
3. Maximal resolved probe (Prop.~\ref{prop:john}) & $n=4$, $9$ walls: $4$ contacts at distance $1.0000$; tight frame to $1.3\times10^{-14}$; $\sum\lambda=4.0000$\\
4. Dilation freedom (Prop.~\ref{prop:dilfree}) & $dg_a-2g_a=0$ exactly for homogeneous walls; $O(1)$ for affine walls\\
5. Conic Steiner formula & layer-1 cell, $n=10$: $3.179\pm0.020$ vs $3.188$; $17.89\pm0.24$ vs $17.93$; $20.12\pm0.16$ vs $20.33$; statistical dimension $3.19$\\
6. Sticky localization (Section~\ref{sec:process}) & $\E\log\det A_T=18.700\pm0.050$ vs It\^o $18.683\pm0.010$; contacts up to $6=\dim\mathrm{Sym}_3$\\\bottomrule
\end{tabular}
\end{center}

\begin{thebibliography}{99}\small
\bibitem{stage13} Stage 13 note: \emph{the leaderboard system read through convex geometry} (this repository, \texttt{notes/stage13}).
\bibitem{stage12} Stage 12 note: \emph{the network as its own two-tier system} (\texttt{notes/stage12}).
\bibitem{stage11} Stage 11 note: \emph{tilings, Bratteli diagrams and singular foliations} (\texttt{notes/stage11}).
\bibitem{Kl} B.~Klartag, Lattice packing of spheres in high dimensions using a stochastically evolving ellipsoid, arXiv:2504.05042.
\bibitem{Sah} J.~Sahasrabudhe, Probabilistic combinatorics at exponentially small scales, arXiv:2512.15077.
\bibitem{Voronoi} G.~Voronoi, Nouvelles applications des param\`etres continus \`a la th\'eorie des formes quadratiques I, J.~reine angew.~Math.~133 (1908) 97--178; J.~Martinet, \emph{Perfect Lattices in Euclidean Spaces}, Springer 2003.
\bibitem{Ball} K.~Ball, Ellipsoids of maximal volume in convex bodies, Geom.~Dedicata 41 (1992) 241--250 (John's decomposition).
\bibitem{CvS} A.~Connes, W.~D.~van Suijlekom, Spectral truncations in noncommutative geometry and operator systems, CMP 383 (2021) 2021--2067.
\bibitem{Takesaki} M.~Takesaki, \emph{Theory of Operator Algebras II}, Springer 2003 (the Plancherel weight of a locally compact group, Ch.~VII); U.~Haagerup, On the dual weights for crossed products of von Neumann algebras I, II, Math.~Scand.~43 (1978).
\bibitem{Bell} J.~Bellissard, Gap labelling theorems for Schr\"odinger operators, in \emph{From Number Theory to Physics}, Springer 1992; J.~Bellissard, D.~Herrmann, M.~Zarrouati, Hulls of aperiodic solids and gap labelling theorems, CRM Monogr.~13 (2000).
\bibitem{Connes} A.~Connes, \emph{Noncommutative Geometry}, Academic Press 1994 (foliation algebras, transverse measures, K-theory of solvable Lie groups via the Thom isomorphism).
\bibitem{BC} J.-B.~Bost, A.~Connes, Hecke algebras, type III factors and phase transitions with spontaneous symmetry breaking in number theory, Selecta Math.~1 (1995) 411--457; A.~Connes, M.~Marcolli, From physics to number theory via noncommutative geometry, 2006.
\bibitem{ALMT} D.~Amelunxen, M.~Lotz, M.~B.~McCoy, J.~A.~Tropp, Living on the edge: phase transitions in convex programs with random data, Inf.~Inference 3 (2014) 224--294; M.~B.~McCoy, J.~A.~Tropp, From Steiner formulas for cones to concentration of intrinsic volumes, DCG 51 (2014) 926--963.
\bibitem{Mahler} K.~Mahler, On lattice points in $n$-dimensional star bodies I, Proc.~R.~Soc.~A 187 (1946) 151--187 (compactness criterion).
\bibitem{Szego} G.~Szeg\H{o}, \emph{Orthogonal Polynomials}, AMS 1939 (Plancherel--Rotach asymptotics, turning points).
\bibitem{Siegel} C.~L.~Siegel, A mean value theorem in geometry of numbers, Ann.~of Math.~46 (1945) 340--347; C.~A.~Rogers, Mean values over the space of lattices, Acta Math.~94 (1955) 249--287.
\end{thebibliography}
\end{document}
